\chapter{Brain Fog}

\section*{Definition}
Brain fog is a colloquial term for subjective cognitive impairment characterized by difficulty concentrating, mental fatigue, forgetfulness, slowed thinking, feeling "fuzzy" or "cloudy," and reduced mental clarity. It is not a formal medical diagnosis but a symptom complex associated with numerous medical, psychiatric, and lifestyle conditions.

\section*{When to See a Doctor}

\begin{itemize}
 \item Cognitive symptoms interfere with work performance, academic function, or daily activities
 \item Brain fog is progressive or accompanied by focal neurological deficits
 \item Associated with severe fatigue, weight changes, heat/cold intolerance (endocrine evaluation)
 \item Onset after starting a new medication
 \item Sudden confusion, trouble speaking, facial droop, one-sided weakness, seizure, or severe headache requires emergency care.
\end{itemize}

\section*{How It Develops}
Brain fog is a nonspecific symptom with multiple possible contributors, including poor sleep, stress, illness, pain, mood disorders, medication effects, and nutritional or endocrine problems. A proposed mechanism in some conditions involves inflammation and altered attention or information processing, but no single biological explanation applies to everyone.

\section*{Causes}
\begin{itemize}
 \item \textbf{Sleep disorders:} Insomnia, sleep apnea, restless legs syndrome, insufficient sleep
 \item \textbf{Autoimmune/inflammatory:} Systemic lupus erythematosus, multiple sclerosis, Sj\"{o}gren syndrome, rheumatoid arthritis, chronic fatigue syndrome
 \item \textbf{Endocrine:} Hypothyroidism, Hashimoto thyroiditis, adrenal insufficiency, diabetes, menopause
 \item \textbf{Nutritional:} Vitamin B12, folate, or iron deficiency; other deficiencies are considered only when risk factors are present
 \item \textbf{Medication-related:} Anticholinergics, benzodiazepines, antihistamines, statins, opioids, chemotherapy (chemo brain)
 \item \textbf{Infectious or post-infectious:} Long COVID and other illnesses; tests for specific infections should be guided by exposure and symptoms
 \item \textbf{Psychiatric:} Depression, anxiety, ADHD, post-traumatic stress disorder
 \item \textbf{Other:} Chronic pain, fibromyalgia, celiac disease, food sensitivities, environmental toxins
\end{itemize}

\section*{Related Symptoms}
\begin{itemize}
 \item Difficulty concentrating, staying on task, or multitasking
 \item Memory lapses (forgetting words, appointments, tasks)
 \item Mental fatigue and slowed processing speed
 \item Feeling of mental "cloudiness" or detachment
 \item Fatigue, sleep disturbance, mood changes
\end{itemize}
\textbf{Important:} Brain fog is a subjective complaint; correlation with objective cognitive testing is often poor, but the symptom significantly impacts quality of life and deserves a thorough evaluation.

\section*{Medical Evaluation}
\begin{itemize}
 \item Self-reported cognitive complaints with history of onset, triggers, and associated symptoms
 \item Cognitive screening: MoCA (Montreal Cognitive Assessment) or Mini-Cog
 \item Laboratory tests are guided by the history and examination and may include blood count, glucose, thyroid testing, vitamin B12, or other tests for suspected reversible causes; broad panels are not routine.
 \item Sleep study (polysomnography) if sleep apnea or other sleep disorder is suspected
 \item Neuropsychological testing for objective quantification of deficits
 \item MRI brain if focal signs, seizures, or progressive cognitive decline are present
\end{itemize}

\section*{Treatment Options}
\begin{itemize}
 \item Treatment of identified underlying cause (e.g., thyroid replacement, B12 supplementation, CPAP for sleep apnea)
 \item Discontinuation or dose reduction of offending medications where possible
 \item Cognitive rehabilitation and compensatory strategies
 \item Exercise prescription and dietary optimization
 \item Stimulants such as modafinil or methylphenidate are not routine treatment for brain fog and should be considered only by a specialist for a defined indication.
 \item Cognitive-behavioral therapy for associated anxiety or depression
\end{itemize}

\section*{Self-Care and Home Management}
\begin{itemize}
 \item Sleep hygiene: consistent schedule, 7--9 hours, limit screen exposure before bed
 \item Regular aerobic exercise (30 minutes, 5 times per week) to improve cerebral blood flow
 \item Mediterranean diet rich in omega-3 fatty acids, antioxidants, and B vitamins
 \item Cognitive strategies: single-tasking, list-making, calendar use, breaks every 60--90 minutes
 \item Mindfulness meditation and stress reduction (reduces cortisol)
 \item Limit alcohol and avoid excessive caffeine; reduce screen time if it disrupts sleep.
\end{itemize}

\section*{References}
\begin{thebibliography}{99}
\bibitem{brain_fog:ref1} Theoharides TC, Conti P. ``Dementia and brain fog: a proposed link between mast cells and microglia.'' \textit{Int J Mol Sci}. 2020;21(6):2043.
\bibitem{brain_fog:ref2} Kao YC, Ho LH, et al. ``Brain fog in long COVID: a review.'' \textit{J Med Virol}. 2023;95(1):e28371.
\bibitem{brain_fog:ref3} Ross AJ, Medow MS, et al. ``Brain fog in patients with postural tachycardia syndrome.'' \textit{Clin Auton Res}. 2013;23(5):275-283.
\bibitem{brain_fog:ref4} Ocon AJ. ``Caught in the thickness of brain fog: exploring the cognitive symptoms of chronic fatigue syndrome.'' \textit{Front Physiol}. 2013;4:63.
\bibitem{brain_fog:ref5} UpToDate. ``Evaluation of cognitive impairment and dementia.'' Wolters Kluwer, 2023.
\end{thebibliography}
